\section{Symmetry}\label{sec:symmetry}

The Symmetry axis collects obstructions caused by involutions,
branch selection, and anomaly-style failures.  In this chapter a
symmetry is not merely a convenient change of notation.  It is part
of the apparatus that controls which branches, traces, and quotient
classes may be distinguished without adding new breaking data.  The
basic question is therefore familiar: when a construction is
invariant under the declared symmetry, what can it select, transport,
or record?  In the running example of \Cref{sec:intro}, the involution
exchanging \(1\) and \(2\) fixes every \(q\)-class; its orbit \(\{1,2\}\) has no fixed point, so by \Fref{F15a} no equivariant section of the orbit map exists (every equivariant map from the orbit set to \(X\) takes values in the fixed points \(3,4\)); choosing \(1\) or \(2\) needs breaking data, such as the readout \(F\), which separates them.

The five laws proceed from fixed data to failure modes and then to
equivalence.  Duality Fixity, or Self-Dual Trace Confinement, is the
layer-agnostic form of the confinement theorem of
\citet{Tsiokos2026SDTC}, which \citet{Tsiokos2026RH} applies to the
zeros of the Riemann zeta function.  For a readout that is strongly
measurable and square-integrable when the visible weight is a measure,
a positive anti-invariant ledger dominated by bounds whose traces have
infimum at most \(0\) is zero, and a separating readout then confines
visible mass to the fixed locus.  Such bounds exist exactly when
\(\psi^-=0\) \(\mu\)-almost everywhere (last clause of \Cref{thm:F14}),
so all content lies in constructing them; the source papers take them
as input, and no result of this paper assumes the Riemann hypothesis.  Selection Obstruction / Repair
Normal Form is a direct theorem: exact symmetry cannot select a
non-fixed branch without declared breaking data.  Spontaneous
Symmetry Breaking / Branch-Selection Formation characterizes when a
branch-selection record exists: exactly when some ground state is not invariant
under the symmetry.  Anomaly / Symmetry Obstruction and Duality
Equivalence are direct theorems about failed symmetry descent and
role-preserving equivalence of formed packages.

The direct selection-obstruction theorem reuses the
Descent--Repair split-pair discipline: a forbidden selection is a
split-pair obstruction at the level of symmetry orbits.  The later
anomaly and duality-equivalence subsections keep the same posture:
the apparatus must say when symmetry descends, when it fails, and
when two symmetry presentations preserve the same formed roles.

\subsection{Duality Fixity (Self-Dual Trace Confinement) [\texorpdfstring{\Fref{F14}}{F14}]}\label{sec:symmetry:duality-fixity}

\begin{definition}[Involutive ledger; anti-invariant readout]\label{def:involutive-ledger}
An involutive ledger is a carrier \(X\) equipped with an involution
\(\Jinv:X\to X\), so \(\Jinv^2=\mathrm{id}\), and a visible weight
\(\mu\ge 0\) on \(X\) (a finite weight function, or a measure).  Its fixed
locus is
\[
  \Fixinv=\{x\in X:\Jinv(x)=x\}.
\]
An anti-invariant readout is a map \(\psi^-\) from \(X\) to a Hilbert space
\(\mathcal H\), with a declared unitary involution \(U\) of \(\mathcal H\), such
that \(\psi^-(\Jinv x)=-U\psi^-(x)\) for every \(x\); when \(\mu\) is a
measure, \(\psi^-\) is strongly (Bochner) measurable, for instance Borel measurable with
\(\mathcal H\) separable, with \(\int_X\|\psi^-\|^2\,d\mu<\infty\).
A map \(\psi^-:X\to\mathcal H\), anti-invariant or not, is \emph{separating} when \(\{x\notin\Fixinv:\psi^-(x)=0\}\) is contained in a
\(\mu\)-null set; for a finite weight this means that
every \(x\) with \(\mu(x)>0\) and \(\psi^-(x)=0\) lies in \(\Fixinv\).
\end{definition}

\paragraph{Reading the statement.}
The objects are those of \citet{Tsiokos2026SDTC}, over a real Hilbert
space.  There the odd part \(\psi^-\) of an observable satisfies
\(\psi^-(\Jinv x)=-\psi^-(x)\), which is the case \(U=\mathrm{id}\) of
\Cref{def:involutive-ledger}; the response involution of that paper,
under which \(\psi^-\) is odd, is a different map and is not the \(U\)
of this definition, so the restriction on \(U\) below is not imposed by
that paper. The anti-invariant ledger
is the positive element
\(A_X=\int_X\psi^-(x)\psi^-(x)^*\,d\mu(x)\), a Bochner integral of
trace-class operators, which converges because
\(\int_X\|\psi^-\|^2\,d\mu<\infty\); for a finite weighted ledger it is the
finite sum \(A_X=\sum_{x}\mu(x)\,\psi^-(x)\psi^-(x)^*\) (Definition~4.1
there). Anti-invariance identifies \(A_X\) as the anti-invariant ledger.
The trace-squeeze argument does not use anti-invariance; the implication from confinement to \(A_X=0\), under the stated restriction on \(U\), does. No clause uses \(\Jinv^2=\mathrm{id}\), nor that \(U\) is unitary or an involution: \(\Jinv\) enters only through \(\Fixinv\), and \(U\) only through \(U\psi^-(x)=-\psi^-(x)\) at fixed points, so every clause holds for an arbitrary map \(\Jinv:X\to X\) and an arbitrary map \(U:\mathcal H\to\mathcal H\). The ledger \(A_X\) lies
in a positive cone with order \(\preceq\) and trace \(\trAX\); positive
trace-class operators on a Hilbert space, or positive semidefinite matrices
with the usual trace, are the classical model. The theorem uses two
properties of the cone and one property of \(A_X\), all of which hold in that
model:
(C1) the trace is monotone, \(A\preceq B\) implies \(\trAX A\le\trAX B\);
(C2) a positive element of trace zero is zero;
(C3) the trace identity \(\trAX A_X=\int_X\|\psi^-(x)\|^2\,d\mu(x)\)
(Lemma~4.3 there; in the classical model it is the trace of a rank-one
operator, \(\trAX(vv^*)=\|v\|^2\)).
The theorem is a reformulation, not a way of establishing
confinement: by its last clause, suitable bounds exist exactly when
\(\psi^-=0\) \(\mu\)-almost everywhere.  Its value lies in isolating the
properties (C1)--(C3) that the squeeze uses.  For a linear involution
\(U\), the condition in the theorem that no nonzero \(v\) satisfies
\(Uv=-v\) holds exactly when \(U=\mathrm{id}\).
The \(B_n\) are the bounds of a domination ladder, a sequence of elements of
the cone with \(A_X\preceq B_n\) for every \(n\); in the source paper they
have the form \(B_n=K^-_n+E_n\) (Definition~6.1 there).

\begin{theorem}[Duality Fixity / Self-Dual Trace Confinement]\label{thm:F14}
Let \(X\) carry a map \(\Jinv:X\to X\) (an involution in the intended reading; \(\Jinv^2=\mathrm{id}\) is not required) and a visible weight \(\mu\) as in \Cref{def:involutive-ledger} (in the intended reading, the ledger of a formed closure; formedness is not used), with
separating readout \(\psi^-:X\to\mathcal H\), strongly measurable and square-integrable when \(\mu\) is a measure, as in \Cref{def:involutive-ledger} but not necessarily anti-invariant, and let its
ledger \(A_X\) (the anti-invariant ledger when \(\psi^-\) is anti-invariant) lie in a positive cone with order \(\preceq\) and trace
\(\trAX\) satisfying (C1)--(C3). If
\[
  A_X \preceq B_n\ \ \text{for every } n
  \qquad\text{and}\qquad
  \inf_n\trAX(B_n)\le 0
  \quad(\text{for every }\varepsilon>0\text{ some }\trAX(B_n)<\varepsilon),
\]
then \(A_X=0\) and \(X\setminus\Fixinv\) is contained in a \(\mu\)-null set,
so \(\mu(X\setminus\Fixinv)=0\) whenever \(\Fixinv\) is measurable (for
instance, when \(\Jinv\) is measurable on a standard Borel space): visible
mass is confined to the fixed locus \(\Fixinv\). Moreover, if \(\psi^-(\Jinv x)=-U\psi^-(x)\) for a declared map \(U:\mathcal H\to\mathcal H\) (a unitary involution in the intended reading), \(\trAX(0)=0\) (true in the classical model), and no nonzero \(v\in\mathcal H\) has \(Uv=-v\) (for instance
\(U=\mathrm{id}\)), then \(A_X=0\) if and only if \(X\setminus\Fixinv\) is
contained in a \(\mu\)-null set; without this condition on \(U\) the
equivalence can fail. Conversely, if \(\psi^-=0\) \(\mu\)-almost everywhere, then \(\trAX A_X=0\) by (C3), and the constant sequence \(B_n:=A_X\) satisfies both displayed hypotheses; hence such a sequence exists if and only if \(\psi^-=0\) \(\mu\)-almost everywhere.
\end{theorem}
\begin{proof}
By (C3), \(\trAX A_X=\int_X\|\psi^-\|^2\,d\mu\ge 0\), and by (C1),
\(\trAX A_X\le\trAX B_n\) for every \(n\). Taking the infimum over \(n\) gives
\(\trAX A_X\le 0\), so \(\trAX A_X=0\). Since \(A_X\) is positive, (C2)
gives \(A_X=0\). By (C3) again, \(\int_X\|\psi^-\|^2\,d\mu=0\), and the
integrand is nonnegative, so \(\psi^-(x)=0\) for \(\mu\)-almost every
\(x\). Then \(X\setminus\Fixinv\) is contained in the union of the null set
\(\{\psi^-\neq0\}\) and the set \(\{x\notin\Fixinv:\psi^-(x)=0\}\), which is
contained in a null set by separation; so \(X\setminus\Fixinv\) is contained
in a \(\mu\)-null set, and \(\mu(X\setminus\Fixinv)=0\) when \(\Fixinv\) is
measurable. For the equivalence under the restriction on \(U\), if \(A_X=0\) then \(\trAX A_X=\trAX(0)=0\), and the argument just given, from \(\trAX A_X=0\) on, confines \(X\setminus\Fixinv\) to a null set. Conversely, let
\(U\) have no nonzero \(v\) with \(Uv=-v\), and let \(X\setminus\Fixinv\) lie in
a null set. For \(x\in\Fixinv\), anti-invariance gives
\(\psi^-(x)=\psi^-(\Jinv x)=-U\psi^-(x)\), so \(U\psi^-(x)=-\psi^-(x)\) and
\(\psi^-(x)=0\); hence \(\psi^-=0\) almost everywhere,
\(\int_X\|\psi^-\|^2\,d\mu=0\), \(\trAX A_X=0\) by (C3), and \(A_X=0\) by (C2).
Without the condition, take one point \(x\) of mass one with \(\Jinv x=x\),
\(Uv=-v\) and \(\psi^-(x)=v\neq0\): then \(X\setminus\Fixinv=\varnothing\) but
\(A_X=vv^*\neq0\). For the converse, if \(\psi^-=0\) \(\mu\)-almost everywhere, (C3) gives \(\trAX A_X=\int_X\|\psi^-\|^2\,d\mu=0\); the constant sequence \(B_n:=A_X\) satisfies \(A_X\preceq B_n\) by reflexivity of \(\preceq\), and \(\trAX(B_n)=0\).
\end{proof}

\paragraph{What the sources supply.}
The theorem assumes the domination ladder; it does not construct one.
\PaperSDTC\ \citep{Tsiokos2026SDTC} proves the same squeeze for real
Hilbert spaces, adds a quantitative form of the confinement and
variants on finite windows, and names as a hypothesis, Self-Dual Trace
Confinement, that in specified self-dual settings such domination
records are supplied by independent structure.  \PaperRH\
\citep{Tsiokos2026RH} applies the squeeze to the reflection
\(s\mapsto1-\bar s\) in the critical line.  Its conditional theorem uses
a named hypothesis on an auxiliary ledger of represented zeros, with
explicit encoding conditions, rather than the global ledger
\(A_Z=\sum_\rho m_\rho(\operatorname{Re}\rho-\tfrac12)^2\) over the
nontrivial zeros, whose vanishing is equivalent to the Riemann
hypothesis for positive integer weights \(m_\rho\).  Its concrete
version uses finite windows of actual zeros that are closed under
\(s\mapsto1-\bar s\), carry unit weights and eventually contain every
nontrivial zero.  The named hypothesis holds exactly when every
represented zero has critical coordinate; it is therefore equivalent to
the Riemann hypothesis on the bare shell, which represents exactly the
nontrivial zeros with coordinate \(\operatorname{Re}\rho\), and the
concrete window version is likewise equivalent to the Riemann
hypothesis.  That paper presents its conditional theorem as a
reformulation, not as evidence.  \Cref{thm:F14} states the
confinement law for any involutive ledger with a separating readout.


\paragraph{Nonclaims.} This subsection does not construct a domination
ladder, does not prove the RH specialization, and does not instantiate the
law to a specific arithmetic carrier.

\subsection*{Layer instantiations}\label{layer-inst:F14}

\begingroup\small
\begin{description}
\item[\emph{Symmetric/skew decomposition under transpose involution
(linear algebra).}] \textit{Analogy.}
For any \(A\in\mathrm{Mat}_n(\mathbb R)\), the transpose
\(J(A)=A^{\!\top}\) is an involution; the unique decomposition
\(A=\tfrac12(A+A^{\!\top})+\tfrac12(A-A^{\!\top})\) separates
\(J\)-symmetric and \(J\)-skew parts.  The skew part \(A-A^{\!\top}\), which \(J\) sends to its negative,
is the \Fref{F14} anti-invariant readout; its Frobenius norm
\(\lVert A-A^{\!\top}\rVert_F\) is \(J\)-invariant, and its vanishing
places \(A\) in the fixed locus \(\mathrm{Sym}_n(\mathbb R)\), on
which the Loewner order is defined
\citep{HornJohnson2013MatrixAnalysis}.

\end{description}
\endgroup


\subsection{Selection Obstruction / Repair Normal Form [\texorpdfstring{\Fref{F15a}}{F15a}]}\label{sec:symmetry:selection-obstruction}

\begin{definition}[Selection obstruction]\label{def:selection-obstruction}
Let a declared symmetry group \(G\) act on a carrier \(X\), and let
\[
  p:X\to X/G
\]
be the orbit quotient.  For a selectedness predicate
\(\chi:X\to\{0,1\}\), the selection obstruction is the split-pair set
\[
  \splitObs(p,\mathrm{id}_X,\chi) =
  \{(x,x')\in X\times X:
      p(x)=p(x')\ \text{and}\ \chi(x)\ne\chi(x')\}.
\]
\end{definition}

\paragraph{Reading the statement.}
The symbols are those of \Cref{thm:F15a} below: \(G\) acts on \(X\)
by \(\rho_X\) and on \(O\) by \(\rho_O\), \(\pi:X\to O\) is
surjective, \(s\) is the selectedness predicate, and
\(\sigma:O\to X\) is a map, in the intended reading a section picking one point in each \(\pi\)-fiber; \(\pi\circ\sigma=\mathrm{id}_O\) is assumed only in the equivariant-section clause.
\(\operatorname{EquivariantSelector}(\rho_O,\rho_X,\sigma)\) says
that this choice commutes with the group:
\(\sigma(\rho_O(g,o))=\rho_X(g,\sigma(o))\) for all \(g,o\).
\(\operatorname{TrivialOrbitAction}(\rho_O)\) says that every group
element fixes every element of \(O\): \(\rho_O(g,o)=o\).
\(\operatorname{HasSelectionSplit}(\pi,s)\) says that some
\(\pi\)-fiber contains a selected and an unselected point.
Under the orbit-quotient hypothesis of the theorem's moreover
clause, these fibers are the orbits, \(O\) plays the role of \(X/G\),
\(\pi\) that of \(p\), and \(s\) that of \(\chi\), so
\(\mathcal O_s\) is the selection obstruction of
\Cref{def:selection-obstruction}.

\begin{theorem}[Selection Obstruction / Repair Normal Form]\label{thm:F15a}
Let \(G\) act on a set \(X\) and on an orbit set \(O\) by
\(\rho_X:G\times X\to X\) and \(\rho_O:G\times O\to O\).  Let
\(\pi:X\to O\) be a surjective map (in the intended reading the orbit map; the
orbit-quotient property is assumed only in the clause that states it), let
\(s:X\to\{0,1\}\) be a selectedness predicate, let
\(\sigma:O\to X\) be a map (in the intended reading a section of \(\pi\), though \(\pi\circ\sigma=\mathrm{id}_O\) is not assumed; the equivariant-section clause below concerns maps that satisfy it).  Define
\[
  \mathcal O_s
  :=
  \{(x,x')\in X^2:\pi(x)=\pi(x')\ \text{and}\ s(x)\ne s(x')\}.
\]
Then
\begin{align*}
  &\operatorname{EquivariantSelector}(\rho_O,\rho_X,\sigma)
    \wedge \operatorname{TrivialOrbitAction}(\rho_O)\\
  &\hspace{4em}\Longrightarrow
    \forall o\in O,\ \forall g\in G,\
      \rho_X(g,\sigma(o))=\sigma(o),\\
  s\ \text{descends through}\ \pi
  &\iff
  \mathcal O_s=\emptyset
  \iff
  \exists\,\bar s:O\to\{0,1\}\ \text{with}\ \bar s\circ\pi=s,\\
  \operatorname{HasSelectionSplit}(\pi,s)
  &\iff
  \mathcal O_s\ne\emptyset.
\end{align*}
If moreover \(\pi\) is the orbit quotient,
\(\pi(x)=\pi(x')\iff x'=\rho_X(g,x)\) for some \(g\in G\), and \(\rho_O\) is
the induced action, \(\rho_O(g,\pi(x))=\pi(\rho_X(g,x))\), then
\(\operatorname{TrivialOrbitAction}(\rho_O)\) holds, every equivariant map
\(\sigma:O\to X\), in particular every equivariant section, takes values in the
common fixed-point set \(X^G\), and if \(X^G=\emptyset\) and
\(O\ne\emptyset\), no equivariant map \(O\to X\) exists; and an equivariant section exists if and only if \(G\) acts trivially on \(X\), and it is then the inverse of \(\pi\).
For every breaker readout \(b:X\to B\), the breaker extension
\(\langle\pi,b\rangle:X\to O\times B\) retains \(\pi\) and carries \(b\), strictly
refines \(\pi\) when \(b\) separates two points of one \(\pi\)-fiber, and factors
through every map \(p\) through which both \(\pi\) and \(b\) factor. Conversely to the fixed-point clause above, under the orbit-quotient and induced-action hypotheses every map \(O\to X^G\) is equivariant, so an equivariant map \(O\to X\) exists if and only if \(X^G\ne\emptyset\) or \(O=\emptyset\). Adjoining \(b\) records the
distinction carried by \(b\); taking \(b=s\) carries selectedness, but the
extension alone does not construct a selector.
\end{theorem}
\begin{proof}
Assume the actions \(\rho_X,\rho_O\), the surjective map
\(\pi:X\to O\), the selectedness predicate \(s\), and the map
\(\sigma:O\to X\).  If \(\sigma\) is equivariant and the \(G\)-action on
\(O\) is trivial, then for every \(g\in G\) and \(o\in O\), equivariance of
the selector and triviality of the action on \(O\) give
\(\rho_X(g,\sigma(o))=\sigma(\rho_O(g,o))=\sigma(o)\), so
\(\sigma(o)\) lies in the common fixed locus of the \(G\)-action on
\(X\).

For selectedness, if \(s=\bar s\circ\pi\), then
\(\pi(x)=\pi(x')\) implies \(s(x)=s(x')\), hence
\(\mathcal O_s=\emptyset\).  Conversely, if
\(\mathcal O_s=\emptyset\), define
\[
  \bar s(\pi(x)):=s(x).
\]
The empty obstruction makes the definition independent of the
representative \(x\), and surjectivity of \(\pi\) supplies
representatives for all \(o\in O\).  If \(\mathcal O_s=\emptyset\), selectedness is a quotient-level
predicate. If \(\mathcal O_s\ne\emptyset\), a split pair shows that no
\(\bar s\) exists; by the breaker clause with \(b=s\), the extension
\(\langle\pi,s\rangle\) removes every split pair, since \(s\) factors through
it.

For the orbit-quotient clause, \(\rho_O(g,\pi(x))=\pi(\rho_X(g,x))=\pi(x)\),
since \(\rho_X(g,x)\) lies in the orbit of \(x\); as \(\pi\) is surjective, the
action on \(O\) is trivial. The first display then places every value of an
equivariant map \(\sigma:O\to X\) in \(X^G\), and if \(X^G=\emptyset\) while \(O\) has an
element \(o\), the value \(\sigma(o)\) cannot exist. For an equivariant section \(\sigma\), the value \(\sigma(\pi x)\) is fixed and lies in the orbit of \(x\), since \(\pi(\sigma(\pi x))=\pi x\); a point of an orbit is fixed only if every point of that orbit is, so \(x\) is fixed and \(\sigma(\pi x)=x\). Conversely, a trivial action makes every orbit a point, so \(\pi\) is bijective and its inverse is an equivariant section. As \(\rho_O\) is trivial, every \(\sigma:O\to X^G\) satisfies \(\sigma(\rho_O(g,o))=\sigma(o)=\rho_X(g,\sigma(o))\), so the equivariant maps are exactly those with values in \(X^G\).

The breaker clauses are the source-repair part of \Cref{thm:F2} with source \(\pi\),
identity transport and target readout \(b\). For the selectedness reading, take
\(b=s\).
\end{proof}

\paragraph{Nonclaims.} The theorem does not classify which
breaking data are allowed and does not instantiate the obstruction to
a particular symmetry group.

\subsection*{Layer instantiations}\label{layer-inst:F15a}

\begingroup\small
\begin{description}
\item[\emph{Vitali sets and the axiom of choice (mathematical
logic).}] \textit{Special case.}
For the free \(\mathbb Q\)-action on \(\mathbb R\) by translation,
no equivariant section of the orbit projection
\(\mathbb R\to\mathbb R/\mathbb Q\) can exist (there are no fixed
points).  The axiom of choice supplies a non-equivariant Vitali
set of representatives, which is necessarily non-measurable: the
\Fref{F15a} declared breaker is the readout \(b(x)=[x=v(\pi(x))]\) for a choice of
representatives \(v:\mathbb R/\mathbb Q\to\mathbb R\), which the axiom of choice supplies, licensing a
non-equivariant selection where no equivariant one is available
\citep{Vitali1905ProblemaMisura}.

\item[\emph{Faddeev--Popov gauge fixing and the Gribov--Singer
obstruction (gauge theory).}] \textit{Analogy.}
Singer's theorem states that the gauge orbit space
\(\mathcal A/\mathcal G\) of Yang--Mills theory admits no
continuous global section.  Faddeev--Popov gauge fixing provides
the \Fref{F15a} declared breaker through a local gauge condition
\(F(A)=0\) together with the Faddeev--Popov determinant
\(\det(\delta F/\delta\omega)\) and ghost fields, licensing local
path-integral computations; the global Gribov--Singer obstruction
witnesses the non-existence of a clean equivariant selector
\citep{FaddeevPopov1967YangMills}.

\item[\emph{Buridan's ass and symmetric tie-breaking (decision
theory).}] \textit{Special case.}
Under exact symmetry between two equally optimal choices, no
equivariant deterministic decision rule selects one option.  Take
\(X=\{\ell,r\}\), \(G=\mathbb Z/2\) exchanging the options, \(O\) a one-point set
and \(\pi\) constant with the induced action; since \(X^G=\varnothing\) and
\(O\ne\varnothing\), \Fref{F15a} gives no equivariant section.  The
declared breaker is a tie-breaking rule --- random coin flip, a
lexicographic preference axis, or a memory of past choices --- that
supplies the non-symmetric apparatus needed for a definite
selection
\citep{Rescher1959ChoiceWithoutPreference}.
\end{description}
\endgroup


\subsection{Spontaneous Symmetry Breaking / Branch-Selection Formation [\texorpdfstring{\Fref{F15b}}{F15b}]}\label{sec:symmetry:spontaneous-ssb}

The setup is a formed closure under an exact declared symmetry, with
the orbit quotient and selection interface already in place as in the
previous subsection.  A spontaneous symmetry-breaking claim asks for
more than a set-theoretic section of the orbit map.  It asks for a
branch-selection record: a map that keeps each state in its orbit and
selects, for some ground state, a branch that the symmetry moves.

\paragraph{Reading the statement.}
A \emph{branch-selection formation record} for \((\mathcal C,G,\Omega)\) is a
triple \((F,c,g)\) with \(F:\mathcal C\to\mathcal C\), \(q\circ F=q\), \(F\)
constant on orbits (\(F(h\cdot x)=F(x)\) for all \(h\in G\), so that
\(F=\sigma\circ q\) for a section \(\sigma\) of \(q\)), \(\Omega(c)\) and
\(g\cdot F(c)\ne F(c)\). The identity map is not such a record, since it is not
constant on a nontrivial orbit: the record selects one branch in each orbit. Clause (o) says that such a record
exists exactly when some ground state is not invariant under \(G\). When
\(\Omega\) is \(G\)-invariant, as in clause (iii), this is the usual definition
of a spontaneously broken symmetry; without that invariance the asymmetry may
already be in \(\Omega\). Clause (i) locates
the selected branch, clause (ii) says that no symmetric selector
produces it, and clause (iii) says that when the ground-state predicate is
symmetric the record exhibits two distinct ground states in one orbit, the
selected branch and its image under \(g\).

\begin{theorem}[Spontaneous Symmetry Breaking / Branch-Selection Formation]\label{thm:F15b}
Let \(\mathcal C\) be a closure space with a \(G\)-action, let
\(\Omega:\mathcal C\to\mathrm{Prop}\) be the ground-state predicate,
and let \(q:\mathcal C\to\mathcal C/G\) be the orbit map onto the
branch quotient.  Then:
\begin{enumerate}[label=(\roman*)]
\item[(o)] a branch-selection formation record exists if and only if some
  ground state is not fixed by \(G\): there are \(c\in\mathcal C\) with
  \(\Omega(c)\) and \(g\in G\) with \(g\cdot c\ne c\);
\end{enumerate}
and for every branch-selection formation record \((F,c,g)\):
\begin{enumerate}[label=(\roman*)]
\item \(F\) maps each orbit into itself, and the branch \(F(c)\) lies in the
  orbit of the ground state \(c\) and is not fixed by \(G\);
\item no \(G\)-equivariant map \(\sigma:\mathcal C/G\to\mathcal C\), for the
  trivial action of \(G\) on \(\mathcal C/G\), takes the value \(F(c)\) at
  \(q(c)\);
\item if \(\Omega\) is \(G\)-invariant, that is \(\Omega(h\cdot x)\iff\Omega(x)\)
  for all \(h\in G\) and \(x\in\mathcal C\), then \(F(c)\) and \(g\cdot F(c)\)
  are distinct ground states in the orbit of \(c\).
\end{enumerate}
So the branch \(F(c)\) cannot be chosen by a symmetric selector: the record
\(F\) is breaker data in the sense of \Cref{thm:F15a}.
\end{theorem}
\begin{proof}
(o) If \(c\) is a ground state and \(g\cdot c\ne c\), choose a section
\(\sigma\) of \(q\) with \(\sigma(q(c))=c\), and let \(F=\sigma\circ q\). Then
\(q\circ F=q\), \(F\) is constant on orbits, and \(F(c)=c\), so
\(g\cdot F(c)\ne F(c)\) and \((F,c,g)\) is a formation record. Conversely, let \((F,c,g)\) be a formation record.
Since \(q(F(c))=q(c)\), there is \(h\in G\) with \(F(c)=h\cdot c\). If \(G\)
fixed \(c\), then \(F(c)=c\) and \(g\cdot F(c)=F(c)\), which is not the case;
so the ground state \(c\) is not fixed by \(G\).

(i) Since \(q\circ F=q\), each \(F(x)\) lies in the orbit of \(x\); in
particular \(F(c)\) lies in the orbit of \(c\), and \(g\cdot F(c)\ne F(c)\)
says that \(G\) does not fix it.  (ii) Let \(\sigma\) be \(G\)-equivariant for
the trivial action on \(\mathcal C/G\).  Then for every \(h\in G\) and every
orbit \(o\), \(h\cdot\sigma(o)=\sigma(h\cdot o)=\sigma(o)\), so every value of
\(\sigma\) is fixed by \(G\), as in the first clause of \Cref{thm:F15a}.  Since
\(g\cdot F(c)\ne F(c)\), \(\sigma(q(c))\ne F(c)\).  The asymmetry is therefore
carried by the selection record \(F\), not by the symmetric quotient.
(iii) By (i), \(F(c)=h\cdot c\) for some \(h\in G\), so \(\Omega(F(c))\) by
invariance, and then \(\Omega(g\cdot F(c))\) by invariance again. Both lie in
the orbit of \(c\), and they are distinct because \(g\cdot F(c)\ne F(c)\).
\end{proof}

\paragraph{Where the breaking comes from.}
The theorem does not produce a non-invariant ground state; whether one exists
is a property of the closure and its ground-state predicate. In a physical
system this is the content of a spontaneous symmetry-breaking result, and the
theorem records where the asymmetry is carried once it holds: in the
selection record, not in the symmetric quotient. \PaperCSL\
\citep{Tsiokos2026CSL} gives a parallel diagnosis in the access setting.


\paragraph{Nonclaims.} This subsection does not exhibit a non-invariant
ground state for any particular closure and does not classify spontaneous
symmetry breaking at a substrate level.  The canonical particle-physics
instances of spontaneous symmetry breaking are surveyed in Goldstone, Salam,
and Weinberg \citep{Goldstone1962BrokenSymmetries}, and the broader
emergent-layer reading of multi-level breaking appears in Anderson's
\emph{More is Different} \citep{Anderson1972MoreIsDifferent}.

\subsection*{Layer instantiations}\label{layer-inst:F15b}

\begingroup\small
\begin{description}
\item[\emph{Bose--Einstein condensation (atomic physics).}] \textit{Analogy.}
When a dilute gas of bosons is cooled below its critical
temperature, a macroscopic fraction of atoms condenses into a
single one-particle quantum state.  The condensate's overall
\(U(1)\) phase is arbitrary --- a continuous symmetry of the
underlying dynamics --- yet each realisation selects a definite
phase, breaking the symmetry.  The 1995 Anderson--Ensher--
Matthews--Wieman--Cornell experiment on rubidium-87 confirmed
Einstein's 1925 prediction; the macroscopic phase selection is the
\Fref{F15b} formation operator
\citep{AndersonEnsherMatthewsWiemanCornell1995BEC}.

\end{description}
\endgroup


\subsection{Anomaly / Symmetry Obstruction [\texorpdfstring{\Fref{F40}}{F40}]}\label{sec:symmetry:anomaly}

Let a formed closure have carrier \(H\), layer quotient
\(q:H\to Q\), declared readouts and audit records, and a would-be
symmetry, an action \(\alpha\) of a group \(G\) on the raw carrier.  The
action is a genuine layer symmetry only when it acts functorially on the quotient and
preserves the declared structure of the layer.

\paragraph{Reading the statement.}
\(\operatorname{SymmetryDescends}(q,\alpha)\) says that each group
element induces a well-defined map on the quotient: for every \(g\),
the map \(h\mapsto q(\alpha(g,h))\) factors through \(q\).
\(\operatorname{DomainDescends}(q,D)\) says that whether \(g\) applies
to \(h\) depends only on the quotient class of \(h\): for every \(g\),
the predicate \(h\mapsto D(g,h)\) factors through \(q\). The other two
conditions are defined by their displayed equations.

\begin{theorem}[Anomaly / Symmetry Obstruction]\label{thm:F40}
Let \(G\) act on a history space \(H\) by
\(\alpha:G\times H\to H\), let \(q:H\to Q\) be a surjective quotient,
let \(\bar\alpha:G\times Q\to Q\) be a candidate descended action,
let \(D:G\times H\to\{0,1\}\) be a domain predicate, and let
\(r:H\to V\) be a readout with value action
\(\rho:G\times V\to V\).  Define
\[
  \mathcal O_{\alpha}
  :=
  \{(g,h,h'):\ q(h)=q(h')\ \text{and}\
    q(\alpha(g,h))\ne q(\alpha(g,h'))\},
\]
and define \(\mathcal O_D\) analogously by replacing
\(q(\alpha(g,h))\) with \(D(g,h)\), and define
\begin{align*}
  \operatorname{GroupLawPreserved}(\bar\alpha)
  &:\Longleftrightarrow
  \forall g,g'\in G,\forall \bar h\in Q,\,
    \bar\alpha(gg',\bar h)=
    \bar\alpha(g,\bar\alpha(g',\bar h)),\\
  \operatorname{ReadoutPreserved}(r,\alpha,\rho)
  &:\Longleftrightarrow
  \forall g\in G,\forall h\in H,\,
    r(\alpha(g,h))=\rho(g,r(h)).
\end{align*}
Then
\begin{align*}
  \operatorname{SymmetryDescends}(q,\alpha)
  &\iff \mathcal O_{\alpha}=\emptyset,\\
  \operatorname{DomainDescends}(q,D)
  &\iff \mathcal O_D=\emptyset.
\end{align*}
Moreover, if \(\bar\alpha\) is a descended action,
\(\bar\alpha(g,q(h))=q(\alpha(g,h))\) for all \(g\) and \(h\), then
\(\mathcal O_\alpha=\emptyset\), \(\operatorname{GroupLawPreserved}(\bar\alpha)\)
holds, and \(\bar\alpha(e,\bar h)=\bar h\) for the identity \(e\) and every
\(\bar h\in Q\), so \(\bar\alpha\) is a group action. Let the would-be symmetry carry a declared record
\(\mathsf{broken}\in\{0,1\}\), which is \(1\) when it is recorded as
explicitly broken. When \(\mathsf{broken}=0\), the \emph{anomaly} is defined by
\[
  \operatorname{Anomaly}
  \;:\iff\;
  \mathcal O_\alpha\ne\emptyset\ \lor\ \mathcal O_D\ne\emptyset\ \lor\
  \neg\operatorname{ReadoutPreserved}(r,\alpha,\rho).
\]
If \(\mathcal O_\alpha\ne\emptyset\), action descent fails and no descended
action exists; when \(\mathsf{broken}=0\) this is an anomaly. If \(\mathcal O_\alpha=\emptyset\), the
\emph{induced action} \(\alpha^{\flat}(g,q(h)):=q(\alpha(g,h))\) is well
defined and is a descended action, hence preserves the group law. An anomaly
then occurs exactly when one of \(\operatorname{SymmetryDescends}(q,\alpha)\),
\(\operatorname{DomainDescends}(q,D)\),
\(\operatorname{GroupLawPreserved}(\alpha^{\flat})\) and
\(\operatorname{ReadoutPreserved}(r,\alpha,\rho)\) fails; since the first and
third hold, this is exactly when domain descent or readout preservation
fails. The record \(\mathsf{broken}\) affects only the anomaly designation:
the descent equivalences and descended-action conclusions hold whatever
its value, while the anomaly characterizations apply only when
\(\mathsf{broken}=0\).
\end{theorem}
\begin{proof}
Assume the action \(\alpha\), quotient \(q\), candidate descended
action \(\bar\alpha\), domain predicate \(D\), readout \(r\), and
value action \(\rho\).  The action descends through \(q\) exactly when
the formula
\[
  \alpha^{\flat}(g,q(h)):=q(\alpha(g,h))
\]
is independent of the representative \(h\).  That independence is
equivalent to \(\mathcal O_{\alpha}=\emptyset\).  The domain predicate
descends by the identical argument applied to \(D(g,h)\), giving
\(\mathcal O_D=\emptyset\).

The displayed equation defines \(\operatorname{GroupLawPreserved}\):
\[
  \bar\alpha(gg',\bar h)
  =
  \bar\alpha(g,\bar\alpha(g',\bar h))
  \quad
  \forall g,g'\in G,\ \bar h\in Q.
\]
If \(\alpha\) is a group action and \(\bar\alpha\) is the descended action,
then for \(\bar h=q(h)\), which covers all of \(Q\) since \(q\) is surjective,
\[
  \bar\alpha(gg',q(h))=q(\alpha(gg',h))=q(\alpha(g,\alpha(g',h)))
  =\bar\alpha(g,q(\alpha(g',h)))=\bar\alpha(g,\bar\alpha(g',q(h))),
\]
so the group law holds. For an action descended from \(\alpha\), surjectivity of
\(q\) also gives \(\bar\alpha(e,q(h))=q(\alpha(e,h))=q(h)\), so the induced map
is a group action. The readout is \(G\)-equivariant exactly when
\[
  r(\alpha(g,h))=\rho(g,r(h))
  \quad
  \forall g\in G,\ h\in H.
\]
If \(\bar\alpha(g,q(h))=q(\alpha(g,h))\) for all \(g,h\) and \(q(h)=q(h')\),
then \(q(\alpha(g,h))=\bar\alpha(g,q h)=\bar\alpha(g,q h')=q(\alpha(g,h'))\),
so \(\mathcal O_\alpha=\emptyset\); conversely a descended action requires
\(\mathcal O_\alpha=\emptyset\). If \(\mathcal O_\alpha=\emptyset\), the induced
action is well defined by the first equivalence and is a descended action, so
it preserves the group law by the computation above; the group law is
therefore not a separate source of anomaly, and an anomaly, as defined in the
statement, is exactly the failure of one of the four conditions for the
induced action. For the induced action, action descent and the group law
hold. In the claimed, unbroken case, an anomaly is therefore exactly a failure
of domain descent or readout preservation. An unclaimed or explicitly broken
symmetry is outside this anomaly designation.
\end{proof}

\paragraph{Nonclaims.} The theorem does not classify anomalies
into carrier-specific forms and does not assert anything about
non-symmetric gates (maps on the carrier, in the sense of
\Cref{sec:apparatus}, that do not commute with the action).  The standard quantum-field-theory treatment of
anomalies, including the Wess--Zumino consistency condition and the
descent-equation classification, is collected in
\citet{Bertlmann2000Anomalies}.

\subsection*{Layer instantiations}\label{layer-inst:F40}

\begingroup\small
\begin{description}
\item[\emph{{IEEE} 754 floating-point non-associativity (numerical
computing).}] \textit{Analogy.}
Associativity of real addition fails in IEEE 754 arithmetic, although addition of finite operands in a fixed format and rounding mode is commutative: \((a+b)+c\) generically differs from \(a+(b+c)\) because each operation rounds. Rebracketing is not a group action; read loosely, with rebracketings of a sum as the would-be symmetry and the rounded value as the readout, the failure is one of ReadoutPreserved; numerical
libraries declare explicit accumulation orders (Kahan or
pairwise summation) as the apparatus extension that mitigates the
anomaly
\citep{IEEE7542019FloatingPoint}.

\end{description}
\endgroup


\subsection{Duality Equivalence [\texorpdfstring{\Fref{F41}}{F41}]}\label{sec:symmetry:duality-equivalence}

Let \(\mathcal L_A\) and \(\mathcal L_B\) be formed closures with
declared symmetry presentations.  Each presentation has an object
quotient, transport data, declared observables, role registry,
residual registry, and status discipline.  A proposed duality is not
just a dictionary between symbols.  It is a pair of translations
between the two formed packages, with the claim that the translations
preserve the roles being compared.

\paragraph{Reading the statement.}
\(\mathcal K_A\) and \(\mathcal K_B\) are the claims of the two
packages, and \(\operatorname{status}_A\), \(\operatorname{status}_B\)
are their declared status assignments, which give each claim its
status in the package's status discipline. The last condition says
that translating a claim carries its status to the corresponding
status of the other package under \(\varsigma\). A \emph{role-preserving
duality equivalence} is defined as the conjunction of the four displayed
conditions; the role readouts live on the quotients, so the role condition
constrains the duality maps themselves. The theorem records this definition
together with its consequences: the inverse map \(\bar G\) also carries routes
to routes and roles back, it is the only two-sided inverse of \(\bar F\), the
role condition fixes \(\rho\) on the image of \(r_A\), and \(F\) and \(G\)
preserve and reflect the quotient classes. These consequences use only the first three conditions together with the descent equations: uniqueness of \(\bar G\) and the fiber clause use the first, the route clause the first and third, and the role clauses the first and second; the status condition, and with it \(\kappa\), \(\varsigma\) and the status maps, enters only the definition.

\begin{theorem}[Duality Equivalence]\label{thm:F41}
Let \(\mathcal C_A,\mathcal C_B\) be sets (in the intended reading, formed closures) with quotients
\(q_A:\mathcal C_A\to Q_A\) and \(q_B:\mathcal C_B\to Q_B\).  Let
\(F:\mathcal C_A\to\mathcal C_B\) and
\(G:\mathcal C_B\to\mathcal C_A\) descend to
\(\bar F:Q_A\to Q_B\) and \(\bar G:Q_B\to Q_A\), meaning explicitly
\(\bar F\circ q_A=q_B\circ F\) and
\(\bar G\circ q_B=q_A\circ G\).  Let
\(r_A:Q_A\to S_A\), \(r_B:Q_B\to S_B\) be role readouts on the quotients,
let \(\rho:S_A\to S_B\) be the role-value map, and let
\(\phi:\Gamma_A\to\Gamma_B\) and
\(\kappa:\mathcal K_A\to\mathcal K_B\) translate routes and claims, let
\(\operatorname{status}_A:\mathcal K_A\to\Sigma_A\) and
\(\operatorname{status}_B:\mathcal K_B\to\Sigma_B\) be the declared
claim statuses, and let \(\varsigma:\Sigma_A\to\Sigma_B\) be the
declared status-value map,
where a route \(\gamma\in\Gamma_A\) acts on \(Q_A\) by
\(\mathrm{act}_A(\gamma):Q_A\to Q_A\) and a route
\(\delta\in\Gamma_B\) acts on \(Q_B\) by
\(\mathrm{act}_B(\delta):Q_B\to Q_B\).
The data form a \emph{role-preserving duality equivalence} when
\begin{align*}
  \bar G\circ\bar F&=\mathrm{id}_{Q_A}
  \quad\text{and}\quad
  \bar F\circ\bar G=\mathrm{id}_{Q_B},\\
  r_B\circ\bar F&=\rho\circ r_A,\\
  \bar F\circ \mathrm{act}_A(\gamma)
  &=\mathrm{act}_B(\phi(\gamma))\circ\bar F
  \quad\forall \gamma\in\Gamma_A,\\
  \operatorname{status}_B(\kappa(k))
  &=\varsigma(\operatorname{status}_A(k))
  \quad\forall k\in\mathcal K_A.
\end{align*}
Then, for a role-preserving duality equivalence, the inverse also
intertwines the routes,
\(\bar G\circ\mathrm{act}_B(\phi(\gamma))=\mathrm{act}_A(\gamma)\circ\bar G\)
for every \(\gamma\in\Gamma_A\), and it transports the roles back:
\(r_B=\rho\circ r_A\circ\bar G\). Moreover \(\bar G\) is the only two-sided inverse of
\(\bar F\), and \(\rho(r_A(a))=r_B(\bar F(a))\) for every \(a\in Q_A\), so if \(r_A\)
is surjective then \(\rho\) is determined by \(\bar F\), \(r_A\) and \(r_B\).
Consequently \(q_A(x)=q_A(x')\iff q_B(Fx)=q_B(Fx')\) for all
\(x,x'\in\mathcal C_A\), and \(q_B(y)=q_B(y')\iff q_A(Gy)=q_A(Gy')\) for all
\(y,y'\in\mathcal C_B\).
\end{theorem}
\begin{proof}
The four displayed conditions define a role-preserving duality equivalence:
the first line makes \(Q_A\) and \(Q_B\) equivalent quotient presentations,
the second transports role observables, the third transports routes, and the
fourth transports claim statuses. We prove the consequences.

For the routes, let \(y\in Q_B\). Using \(\bar F\bar G=\mathrm{id}\),
then the route condition, then \(\bar G\bar F=\mathrm{id}\):
\[
  \bar G\,\mathrm{act}_B(\phi\gamma)\,y
  =\bar G\,\mathrm{act}_B(\phi\gamma)\,\bar F\bar G\,y
  =\bar G\bar F\,\mathrm{act}_A(\gamma)\,\bar G\,y
  =\mathrm{act}_A(\gamma)\,\bar G\,y .
\]
For the roles, \(r_B=r_B\circ\bar F\circ\bar G=\rho\circ r_A\circ\bar G\).
If \(\bar G'\) is another two-sided inverse of \(\bar F\), then \(\bar G'=\bar G'\circ\bar F\circ\bar G=\bar G\). The identity \(\rho(r_A(a))=r_B(\bar F(a))\) is the role condition read pointwise, and it fixes \(\rho\) on the image of \(r_A\).
For the fiber clause: since \(q_B\circ F=\bar F\circ q_A\) and \(\bar F\) is a
bijection, \(q_B(Fx)=q_B(Fx')\) exactly when \(q_A(x)=q_A(x')\); likewise for
\(G\).
\end{proof}

\paragraph{Nonclaims.} The theorem does not classify
symmetry presentations and does not assert a unique role presentation
beyond observable equivalence.

\subsection*{Layer instantiations}\label{layer-inst:F41}

\begingroup\small
\begin{description}
\item[\emph{Pontryagin duality (harmonic analysis).}] \textit{Analogy.}
For a locally compact abelian group \(G\), the Pontryagin dual
\(\widehat G\) of continuous characters \(G\to S^1\) satisfies the
Pontryagin duality theorem
\(G\cong\widehat{\widehat G}\), with the Fourier transform
\(L^2(G)\to L^2(\widehat G)\) a unitary isomorphism.  Pontryagin dualization is a contravariant equivalence: for closed \(H\), the
dual of \(G/H\) is the annihilator subgroup \(H^\perp\subseteq\widehat G\), and
double dualization recovers \(G\). For functions with defined convolution,
Fourier transform exchanges convolution and pointwise multiplication; this is
an analogy with reversible translation, with quotient variance and route data
specified separately \citep{Pontryagin1966TopologicalGroups}.

\end{description}
\endgroup


\subsection*{Symmetry at a glance}\label{sec:symmetry:summary}\addcontentsline{toc}{subsection}{Symmetry at a glance}
\begin{table}[H]
  \centering
  \footnotesize
  \begin{tabularx}{\textwidth}{@{}
      >{\raggedright\arraybackslash}p{0.24\textwidth}
      >{\raggedright\arraybackslash}X
      >{\raggedright\arraybackslash}p{0.12\textwidth}@{}}
    \toprule
    Name & One-line claim & Substrate cite \\
    \midrule
    Duality-Fixity Theorem (\Fref{F14}) &
    For a readout that is strongly measurable and square-integrable when the weight is a measure, and a positive cone satisfying (C1)--(C3) (monotone trace, positive trace-zero elements vanish, trace identity), domination by bounds \(B_n\) with \(\inf_n\operatorname{tr}B_n\le0\) forces the anti-invariant
    ledger to vanish and, for a separating readout, confines visible mass
    to \(\Fixinv\) up to a null set; such bounds exist exactly when \(\psi^-=0\) \(\mu\)-almost everywhere.  &
    \PaperSDTC, \PaperRH \\
    Selection-Obstruction Theorem (\Fref{F15a}) &
    For the orbit quotient with its induced action, an equivariant selector
    takes only fixed points; selectedness descends through the quotient exactly
    when its obstruction is empty.  &
    \textemdash \\
    Spontaneous Symmetry Breaking (\Fref{F15b}) &
    A branch-selection record exists exactly when some ground state is
    not invariant; no symmetric selector produces its branch.  &
    \PaperCSL\ (parallel) \\
    Anomaly Theorem (\Fref{F40}) &
    For a claimed symmetry not recorded as explicitly broken, an anomaly
    is the failure of action descent, domain descent or readout
    preservation; for a surjective layer quotient, a descended action is a
    group action (it satisfies the group law and the identity acts
    trivially).  &
    \textemdash \\
    Duality-Equivalence Theorem (\Fref{F41}) &
    In a role-preserving duality equivalence, the inverse also
    intertwines the routes and \(r_B=\rho\circ r_A\circ\bar G\).  &
    \textemdash \\
    \bottomrule
  \end{tabularx}
  \caption{Symmetry at a glance: the five laws of \Cref{sec:symmetry}. Duality Fixity has restricted alignment.}
  \label{tab:axis-symmetry-summary}
\end{table}

\FloatBarrier
